\chapter{Kerangka Adopsi Cloud dan Strategi Migrasi 7R}

\section*{Tujuan Pembelajaran}
\begin{itemize}
  \item Membandingkan struktur kerangka adopsi \textit{cloud} dari AWS (Amazon Web Services), Microsoft, dan Google.
  \item Menilai kesiapan organisasi dengan enam perspektif AWS CAF.
  \item Memberi usulan strategi 7R yang beralasan untuk setiap aplikasi dalam portofolio.
\end{itemize}

\section{Pendahuluan}

% TODO
% Cakupan: Medikara Health memiliki puluhan aplikasi hasil akuisisi dan tidak tahu harus mulai dari mana; migrasi yang hanya digerakkan divisi TI sering gagal karena orang, proses, dan tata kelola belum siap.

\section{Tiga Kerangka Adopsi Cloud}

% TODO
% Cakupan: AWS CAF (enam perspektif, domain transformasi, dan capability dasar); Microsoft CAF versi 2026 (Strategy, Plan, Ready, Adopt, Govern, Secure, Manage, serta skenario sovereignty); Google Cloud Adoption Framework (Learn, Lead, Scale, Secure dan tiga fase kematangan); persamaan dan perbedaan ketiganya.

\section{Penilaian Kesiapan dengan Enam Perspektif}

% TODO
% Cakupan: perspektif Business, People, Governance, Platform, Security, dan Operations; pertanyaan kunci tiap perspektif; pemangku kepentingan yang perlu dilibatkan; cara memberi skor kesiapan tanpa terjebak formalitas.
% Kasus Medikara: kesiapan tim TI yang tersebar di lima rumah sakit dan dampak temuan audit kontrol akses pada perspektif Security.

\section{Inventaris Portofolio Aplikasi dan Dependensi}

% TODO
% Cakupan: atribut minimal inventaris (pemilik, pengguna, kritikalitas, teknologi, kontrak, dependensi); cara menemukan dependensi tersembunyi; data survei industri bahwa pemahaman dependensi aplikasi adalah hambatan migrasi utama; prinsip arsitektur dan gap analysis disinggung singkat (dibahas rinci di IT40513).
% Pendalaman teknis: memakai data log jaringan atau alat discovery untuk membangun peta dependensi.

\section{Strategi Migrasi 7R}

% TODO
% Cakupan: Retire, Retain, Rehost, Relocate, Repurchase, Replatform, dan Refactor; kasus pemakaian tiap strategi menurut AWS Prescriptive Guidance; alasan Refactor tidak dianjurkan untuk migrasi berskala besar; urutan "pindah dulu, modernisasi kemudian".
% Kasus Medikara: usulan 7R untuk SIMRS, ERP keuangan, Medikara+, sistem laboratorium, dan aplikasi apotek.

\section{Landing Zone dan Cloud Center of Excellence}

% TODO
% Cakupan: landing zone sebagai fondasi akun, jaringan, identitas, dan kebijakan; CCoE sebagai tim lintas fungsi yang memandu adopsi; data survei industri tentang tingkat adopsi CCoE; guardrail tata kelola dibahas lanjut di Bab 14.

\section{Aktivitas: Penilaian Kesiapan dan Portofolio 7R}

% TODO
% Tujuan aktivitas: menghasilkan penilaian kesiapan cloud enam perspektif dan portofolio aplikasi dengan usulan 7R yang beralasan.
% Peran dalam tim: anggota non-teknis memimpin perspektif Business, People, dan Governance; anggota teknis memimpin perspektif Platform, Security, dan Operations; seluruh anggota membahas usulan 7R bersama.
% Langkah: (1) isi lembar penilaian kesiapan; (2) inventaris aplikasi beserta pemilik, pengguna, dependensi, dan kondisinya; (3) beri usulan 7R awal dengan alasan satu kalimat; (4) tandai aplikasi yang usulannya masih diperdebatkan.
% Luaran dan tonggak capstone: C2, penilaian kesiapan cloud dan portofolio aplikasi dengan usulan 7R.
% Berkas sources yang direncanakan: sources/chapter04/lembar-penilaian-kesiapan-caf.xlsx, sources/chapter04/templat-inventaris-aplikasi.xlsx, sources/chapter04/panduan-keputusan-7r.md

\section{Ringkasan}

% TODO

\section{Latihan}

\begin{enumerate}
  \item Jelaskan mengapa panduan AWS menyarankan agar strategi Refactor tidak dipakai untuk sebagian besar aplikasi dalam migrasi berskala besar.
  \item Medikara Health memiliki SIMRS dengan versi berbeda di lima rumah sakit. Usulkan strategi 7R untuk SIMRS, lalu pertahankan usulan tersebut terhadap satu alternatif lain.
  \item Perspektif CAF mana yang paling sering diabaikan ketika migrasi hanya digerakkan oleh divisi TI? Jelaskan akibatnya.
\end{enumerate}
